\section{Several marked insertions on the restarted history}
\label{sec:allarity-planted-compiler}

We keep the actual restarted density from
Section~\ref{sec:restarted-rank-history}. Its additive remainder is
removed only in the auxiliary unmarked normalization. The marked integral
itself must retain that remainder. We first prove the finite algebra that
identifies all the selected insertion coefficients, then explain this
mixed normalization and the extra local price it requires.

\subsection{Selecting a cell in the marker algebra}

For a finite set $V$, let
\[
 \mathscr A_V=\mathbb C[z_v:v\in V]/(z_v^2:v\in V).
\]
Its basis consists of $z_A=\prod_{v\in A}z_v$, $A\subseteq V$.
Products satisfy $z_Az_B=z_{A\cup B}$ for disjoint sets and are zero
otherwise. If $F$ has zero constant coefficient, define
$\exp_*F=\sum_{j=0}^{|V|}F^j/j!$. Every product of more than $|V|$
nonconstant monomials has a repeated variable, so it vanishes. Hence
this finite exponential satisfies the usual addition formula for
commuting elements of zero constant coefficient, by the binomial theorem.
In particular, for every finite family of pieces,
\begin{equation}\label{eq:planted-exp}
 \prod_\alpha(1+y_\alpha z_{A_\alpha})
       =\exp_*\left(\sum_\alpha y_\alpha z_{A_\alpha}\right)
       \quad(A_\alpha\ne\varnothing).
\end{equation}
There is no disjointness assumption on the $A_\alpha$: their overlapping
products vanish on both sides.

For $T\subseteq V$, define $\pi_T$ by setting $z_v=0$ for $v\notin T$.
Thus $\pi_T(\sum_A c_Az_A)=\sum_{A\subseteq T}c_Az_A$.
It is an algebra homomorphism, and therefore commutes with $\exp_*$.
For a finite cell set $\mathcal C$ and mark set $J$, put
$V=\mathcal C\times J$. The cell embedding
$\iota_c:\mathscr A_J\to\mathscr A_V$ sends $z_j$ to $z_{(c,j)}$;
it too is an algebra homomorphism.

Fix one coarse cell $c$ and a selected set $T\subseteq\{c\}\times J$.
The incoming insertion pieces are indexed by pairs $(c',\alpha)$.
Each piece has a nonempty marker support $A_{c',\alpha}$ lying in
$\{c'\}\times J$. Its value in a particular sector, forest, tag and
Gaussian coordinate is $y_{c',\alpha}$. The spatial geometry selects
exactly those pieces for which
$A_{c,\alpha}\subseteq T$; write $\mathcal I_T$ for this selected
index set and write $B_\alpha\subseteq J$ for their marker supports
with the cell coordinate removed. Define the planted polynomial
\[
             P_T=\sum_{\alpha\in\mathcal I_T}
                         y_{c,\alpha}z_{B_\alpha}.
\]
Applying $\pi_T$ to \eqref{eq:planted-exp} proves
\begin{equation}\label{eq:planted-selection}
 \pi_T\prod_{c',\alpha}(1+y_{c',\alpha}z_{A_{c',\alpha}})
           =\pi_T\iota_c(\exp_*P_T).
\end{equation}
To check the selection explicitly, a piece from $c'\ne c$ has at least
one marker outside $T$ and is killed. A piece from $c$ is retained
exactly when its entire marker support is contained in $T$. A discarded
monomial maps to zero, so its factor $1+y_{c',\alpha}z_{A_{c',\alpha}}$
maps to one. In the expanded product, every term selecting a discarded
monomial vanishes. This proves the identity at all orders simultaneously.

The spatially planted value is the same masked field value as in the
left side. Write it as $\Gamma_\alpha(\xi)h_\alpha(U(\xi))$, where
$0\le\Gamma_\alpha\le1$ is its local cutoff. On the region where
$\Gamma_\alpha>0$, the chart's healthy-field condition makes the original
mask equal to one on the full inspected support. Where
$\Gamma_\alpha=0$, both products are zero. Consequently this formula
agrees everywhere with
$\Gamma_\alpha\operatorname{masked}(h_\alpha)(U(\xi))$.
The sector and tag are fixed during this verification. They are not
chosen separately for the different readers.

All linear maps on $\mathscr A_V$ are continuous because this space is
finite dimensional. Hence an integrable polynomial-valued function can
be integrated coefficientwise, and its integral commutes with $\pi_T$
and $\iota_c$. If $\nu_t$ are the finitely many sector/tag measures and
$\lambda_t$ their scalar weights, \eqref{eq:planted-selection} gives
\begin{equation}\label{eq:planted-integral}
 \sum_t\lambda_t\int
       \pi_T\prod_{c',\alpha}(1+y_{c',\alpha,t}z_{A_{c',\alpha}})
          \,d\nu_t
 =\pi_T\iota_c\left(\sum_t\lambda_t\int\exp_*P_{T,t}\,d\nu_t\right).
\end{equation}
Here integrability is provided by the common positive marker weights
and the paid product envelopes: the absolute value of each coefficient
is bounded by its weighted polynomial norm divided by its positive
monomial weight. The Gaussian/Haar envelope bounds that norm by an
integrable function. For the countable activity menu, first apply the
identity to finite truncations. The summable full-support majorants and
coefficientwise countable regrouping of
Section~\ref{sec:ex-lattice-rg-multisource} provide the limiting
integrable coefficients; the finite-dimensional linear maps commute with
those limits. Thus no interchange of a merely formal infinite series is
used in \eqref{eq:planted-integral}.

\subsection{The marked density and its three kinds of factors}

Let $P$ denote the fine representation and $Q$ the coarse one. Write
$\mathcal K$ for a finite family of lifted background packets, with
full spatial supports $X_b$ and fine flag inventories $F_b$. Its
unmarked coefficients reproduce the activities of $P$ at each activity
address. This addresswise identity is stronger than equality of their
summed gases. For regular insertion pieces $h_\alpha$ with nonempty
marker supports, the exact polynomial density is
\begin{equation}\label{eq:planted-density}
 D_P(U)=e^{\operatorname{exponent}_P(U)}
    \exp_*\left(\sum_\alpha z_{A_\alpha}
                    \operatorname{masked}(h_\alpha)(U)\right)
    \sum_{B\ \text{admissible}}\prod_{b\in B}K_b(U).
\end{equation}
Admissibility requires pairwise disjoint full supports and the exact
identity $\bigcup_{b\in B}F_b=\operatorname{flags}_P(U)$.
It is not enough to require just disjoint inventories.

Fix the actual coarse configuration, positioned chart, forest, tag and
Gaussian coordinate. The unmarked chart expansion, including its Haar
Jacobian and Gaussian normalization, has the following sliced form.
For each fine flag set $F$,
\begin{equation}\label{eq:planted-body}
 \mathbf1_{\{\operatorname{flags}_P(U)=F\}}
          J(U)e^{\operatorname{exponent}_P(U)}
   =\Lambda\sum_{A\ \text{compatible}(F)}\prod_{a\in A}p_a.
\end{equation}
Here $J$ denotes the full scalar chart factor, and $\Lambda$ is the
same positive loading used by the unmarked construction. Each $p_a$
has a fixed coarse support, a fine inventory and a coarse inventory.
Compatibility means disjoint coarse supports, fine inventory union $F$,
coarse inventory union equal to the coarse flags, and the sector key
specified by those same packets. These are precisely the data of the
unmarked sector expansion of Section~\ref{sec:gaussian-reassembly}.

To obtain the marked identity, first fix a background packet family
$B$. If its full supports overlap its contribution is zero. Otherwise
multiply \eqref{eq:planted-body}, with $F=\bigcup_{b\in B}F_b$, by
$\prod_{b\in B}K_b$ and by the insertion product
\eqref{eq:planted-exp}. The indicator on the left imposes exactly the
background flag test in \eqref{eq:planted-density}. Expand the insertion
product as a sum over its subsets. On the right we now have three
independent kinds of factors: scalar body packets $p_a$, background
packets $K_b$, and marked insertion pieces. Index their union by the
disjoint union of these three alphabets. This is a bijection between
the three subset choices and the subsets of the union alphabet.
The flag and sector tests are unchanged. If a marker is used twice,
the corresponding product is zero in $\mathscr A_V$, exactly as on
the expanded left side. Summing over $B$ proves the full marked
chart identity at every coefficient.

Integrate with the common Gaussian/Haar product envelopes, sum the
sector partition and average the tags. The coefficientwise argument of
Section~\ref{sec:ex-lattice-rg-multisource} then applies to this same
three-alphabet expansion. It gives the pushed density
$\mathcal T D_P$ represented by the resulting coarse polynomial
activities and their regular exponential. In particular this step does
not follow by repeating the one-address affine argument of
Section~\ref{sec:marked-block-integral}; that argument alone has no
mixed higher coefficients.

\subsection{Restoring the additive remainder before normalization}

Put $P^0=\operatorname{white}(P)$, so that its additive remainder has
been deleted while its coupling, path, masks and forced activities are
unchanged. In the notation of Section~\ref{sec:restarted-rank-history},
$P=e^{-E}P^0$. On the actual chart the additional exponent is a sum
$\sum_i f_i=-E(U)$ over local coarse supports. Suppose its support
prices $b_i\ge0$ satisfy $|f_i|\le b_i$ and have anchored weighted
sum at most $d\le1$. The functions $f_i$ are real-valued. Let $M_i$ be the real-valued
unperturbed local exponents,
with
\[
 |e^{M_i}-1|\le w_i,\qquad w_i\ge0,
 \qquad\hbox{anchored weighted sum of }w_i\le e\le1.
\]
The weights are at least one, so individually $b_i\le d$ and
$w_i\le e$. Hence $e^{M_i}\le2$ and
$|e^{f_i}-1|\le2|f_i|\le2b_i$. The latter elementary inequality
holds for real $|f_i|\le1$, by convexity of the exponential on
$[0,1]$ and $e-1<2$, and by $1-e^{-t}\le t$ for $t\ge0$.
It follows that
\[
 |e^{M_i+f_i}-1|
 \le |e^{M_i}-1|+e^{M_i}|e^{f_i}-1|
 \le w_i+4b_i.
\]
Thus the perturbed auxiliary price is at most $e+4d$. The base Gaussian
quadratic form and the scalar loading have not changed. This proves
\eqref{eq:planted-body} for $P$ itself from the neutral chart expansion:
add $f_i$ to the same local exponents and use their exact sum $-E(U)$.
The full-support locality of these terms was established in the
restarted-history construction, including the final flag tests.

For the actual scale purchases, write $\omega$ for the permitted
perturbation price. The neutral auxiliary price is at most $\omega$
and $d\le\omega$, so the preceding bound is at most $5\omega$ and
in particular at most the reserved $6\omega$. The two other elementary
letter families have anchored prices at most $u^{-1/2}$ each, where
$u$ is the retimed coupling scale. Since $u^{-1/2}\le\omega$, their
combined price is at most $8\omega$. The background large-field
letters retain the decay $\exp(-\alpha\,\mathrm{hot}^2\,\mathrm{shares})$;
the separate purchase $2a_{\rm coast}\le\alpha/4$ pays their forced
inflation. With the reserved line multiplier $L_0$, the integrated
cluster bound is therefore purchased through
\[
                 L_0(e^{8\omega}-1)\le\delta_{\rm step}.
\]
This is the local unmarked price used by the restarted recurrence.
The fixed stage constants and Gaussian envelopes in this calculation
are those of the neutral compiler at the same physical parameters.
The perturbation estimate does not assert derivative bounds for $E$.

There are now two distinct, compatible uses of the representations.
The body identity \eqref{eq:planted-body} is for $P$, including $E$.
The unmarked auxiliary alphabet consists of these \emph{same corrected}
body packets $p_a$, together with the archived background packets of $P^0$.
It does not use the former neutral body packets. Its loading is the same
$\Lambda$ as in \eqref{eq:planted-body}.

The normalization needed here has two precise parts. Write $Z_j$ for
its zero-marker regular components. For every deleted spatial support
$D$ and every coarse configuration $V$, require the identity
\begin{equation}\label{eq:planted-holes}
 \Lambda(V)\,\mathcal G_{Z;D}(V)
 =e^{\operatorname{exponent}_Q(V)}
        \prod_a(1+p_a^{\rm pay}(V))
        \prod_{b:\,X_b^{\rm hole}\cap D\ne\varnothing}
                    (1+h_b(V)).
\end{equation}
Here $\mathcal G_{Z;D}$ is the compatible gas of the $Z_j$ whose
full supports are disjoint from $D$. The factors $p_a^{\rm pay}$ and
$h_b$ are the fixed boundary and hole corrections of the unmarked
regrouping, with their original masks. Secondly, the unmarked coarse
forced family must embed into the forced components of this alphabet,
with identical values, supports and inventories, and must exhaust its
zero-marker forced components. The forced family, together with these
same hole and boundary correction factors, must reproduce the archived
activity of $Q$ at each activity address, with its full support and
flag inventory. This addresswise identity is required before applying
the exclusion regrouping. These are the all-deleted-support
identity and the forced-family embedding supplied by the unmarked
forest and exclusion construction. They must be established for the
corrected $p_a$ and this actual coarse $Q$; equality of the two archived
fine backgrounds alone does not establish them.

We verify that these two statements, once available, pass to the full
marked alphabet. First coalesce background packets with the same
address and marker mode. Their zero-marker values are exactly the
archived activities at that address, by the addresswise background
identity. This embeds the entire unmarked alphabet into the full
alphabet, preserving values, full supports, inventories and sector
keys. Every remaining elementary insertion has a nonempty marker
support. If a full connected component has empty marker support,
every elementary factor in it must therefore have empty marker
support: its marker support is their union. Hence it comes from the
unmarked alphabet. Conversely every unmarked component has the same
connectedness and eligibility tests after the embedding. This is a
bijection of zero-marker components preserving their full supports.
It restricts to the components disjoint from any given $D$. Thus
$\mathcal G_{Z;D}$ is unchanged, and \eqref{eq:planted-holes} holds
for the full alphabet at every deleted support. The same component
bijection transports the forced-family embedding and its exhaustion
property.

Finally fix a compatible marked forced family and let $D$ be the union
of its full supports. Its residual zero-marker gas is precisely
$\mathcal G_{Z;D}$. Substitute \eqref{eq:planted-holes}. A hole factor
occurs exactly when its support meets at least one member of that forced
family, since $D$ is their union. Expanding these factors and regrouping
by their incident forced components is the weight-preserving exclusion
bijection of Section~\ref{sec:forced-exclusion}. The resulting polynomial
has exponent $\operatorname{exponent}_Q$ and the actual coarse archived
background, by the transported forced-family embedding. Summing over the
marked forced families therefore represents $\mathcal T D_P$ with this
coarse vacuum normalization. This uses \eqref{eq:planted-body} for $P$
and \eqref{eq:planted-holes} for the corrected auxiliary alphabet; it
never substitutes $\mathcal T D_{P^0}$ for the marked integral.

This proves the finite algebra and the additive-remainder correction
needed by the all-order compiler. Its quantitative marked use still
requires the common Gaussian reader estimates and their cumulant
approximation, the day and night marker prices, and the selected
spatial support bounds. In particular the cell selection
\eqref{eq:planted-integral} must be applied with the same readers and
same sector measures as the body expansion, at the purchased slice
radius. Those analytic and geometric estimates are separate from the
polynomial identities proved here.
